% Guidance for the Junior Level 3 show-me checklists and the week-1 warm-up
% (input by answer-key.tex; the pages the parent fills in are in assessments.tex)
\kcpart{A}{Assessments}{How to use the show-me checklists}

\begin{teachernote}[Watching, not testing]
\begin{itemize}[itemsep=1pt]
  \item \textbf{Nothing here is a grade.} The checklists tell you what to practice next and let you
        and your child see progress. Share them as \enquote{you can now\dots} and \enquote{next we will\dots}.
  \item \textbf{Watch during normal play.} Most rows show up on their own during the play day (week 5)
        or the project and showcase (weeks 9--10). Keep the sheet on a clipboard and check boxes as you see them.
  \item \textbf{Ask, do not quiz.} If a row has not come up, ask near the end:
        \enquote{Can you show me how you\dots?} Never time it, and never ask in front of
        a group if a child is shy: ask while you sit beside them.
  \item \textbf{Count what they do, not how they say it.} Pointing, acting it out and short answers all count.
        A child who types slowly or spells with help can still check \enquote{Got it} for a Python row:
        the row is about the idea, not typing speed.
  \item \textbf{One box per row,} the best you saw that day. If you are unsure between two columns,
        choose the lower one and look again next week.
\end{itemize}
\end{teachernote}

%% one row of guidance: \smguide{row}{not yet}{getting there}{got it}{what to do next}
\newcommand{\smguide}[5]{%
  \par\needspace{7\baselineskip}\noindent\textbf{\color{kcaccent}#1}\par
  \begin{itemize}[leftmargin=1.4em,itemsep=1pt,topsep=2pt]
    \item[\faSeedling] \textbf{Not yet:} #2
    \item[\faShoePrints] \textbf{Getting there:} #3
    \item[\faCheckCircle] \textbf{Got it:} #4
    \item[\faArrowRight] \textbf{If not yet:} #5
  \end{itemize}\medskip}

%% ---------------------------------------------------------------------
\section{Week 1: the \enquote{show me} warm-up}
The warm-up is not graded and has no score. It tells you what carried over from Junior Level 2,
often after a summer or a long break. Use it like this:
\begin{itemize}[itemsep=2pt]
  \item \textbf{Two or three rows per child are enough.} Choose the pen square (it comes back this
        week) and one or two Level 2 game rows. Check the column you see; leave the rest blank.
  \item \textbf{\enquote{Not yet} on the pen square or the My Block:} fine, week 1 revisits both. Pair
        that child as Helper first with a confident partner, and use the support path of week 1
        (start from a square that is already built).
  \item \textbf{\enquote{Not yet} on forever, if touching or the score:} these come back in the week-5
        play day and in a Scratch Big Build. Make a note, and show the Level 2 block once in free play.
  \item \textbf{The circled pictures} tell you which examples and free-play ideas to offer: a child who
        circles games is likely to choose a Scratch game in week 9, one who circles art or code a Python
        postcard. Mention it again when they plan their Big Build.
  \item \textbf{The faces.} \enquote{Worried} or \enquote{not sure} about writing code with words: say that
        every Python program in this level already works and that they will change it, not start from
        nothing. Check in quietly again in week 6. \enquote{Super excited}: give them the Checker or Helper
        job in week 6, so the excitement turns into reading the code closely.
\end{itemize}
Keep the page in your folder. It has no personal information beyond a nickname.

%% ---------------------------------------------------------------------
\section{Week 5 show-me checklist}

\smguide{Uses a repeat inside a repeat to make a pattern (J3.G1)}
  {builds one square, or puts the second repeat under the first instead of inside it.}
  {builds the nested repeat with a reminder (\enquote{Which repeat goes inside which?}) or forgets the turn between the shapes.}
  {puts a square (repeat 4) inside an outer repeat with a turn, and changes a number to get a new pattern on purpose.}
  {Do the week-1 clap pattern again: repeat 3 of (clap, clap, stomp), with the inner card held inside the outer card. Then build a nested repeat together with you as Helper.}

\smguide{Plans a program in parts before building it (J3.G2)}
  {starts dragging blocks without a plan, or the plan is one big step.}
  {names the parts with help, then builds them, sometimes all at once.}
  {names two or three parts on paper (for example \enquote{start}, \enquote{dance}, \enquote{end}), builds and tests one part at a time.}
  {Plan something non-computer first (the week-2 party): three named parts, one job each. Then plan a tiny program the same way and build only the first part.}

\smguide{Tallies data and makes a bar graph (J3.G3)}
  {counts the answers but loses track, or draws bars that do not match the count.}
  {tallies in groups of five with a reminder, and draws the bars with help to line them up.}
  {tallies the answers, draws one bar per answer at the right height, and says what the graph shows (\enquote{Most people like bananas}).}
  {Do a 3-question survey of the family with only two choices. Use sticky notes as bars, one note per answer, then count them together.}

\smguide{Explains that a computer learns from examples and can learn unfair mistakes (J3.G4)}
  {says the computer \enquote{already knows}, or cannot say how it learned.}
  {says it learned from examples, but not what goes wrong with bad examples.}
  {says, in their own words, that a computer learns from examples and that too few or one-sided examples make it guess wrong, or unfairly.}
  {Play \enquote{teach the robot} from week 4 again with only cats as examples, then show a dog. Ask: \enquote{Why did the robot get it wrong? Which examples would help?}}

\smguide{Matches Scratch blocks to Python words (J3.G5)}
  {cannot find any Python twin yet, or matches by color or length of the card.}
  {matches move and turn, and needs help with repeat or pen color.}
  {matches move with \texttt{t.forward}, turn with \texttt{t.left}, repeat with \texttt{for}, and pen color with \texttt{t.color}, and can say why.}
  {Lay out only three pairs (move, turn, repeat). Read each Python card aloud, act it out (step forward, turn), then find its block.}

\smguide{Finds and fixes a bug (J3.G6)}
  {says \enquote{it's broken} and waits for you to fix it.}
  {finds the bug when you point to the part of the script, then fixes it.}
  {tests it, looks closely, changes one thing and tests again, and can say what was wrong.}
  {Use the steps on the cheat sheet out loud together: test, look, change one thing, test again. Put a working script beside the broken one to compare.}

%% ---------------------------------------------------------------------
\section{Week 10 show-me checklist}

\smguide{Runs a Python program in Thonny (J3.G5)}
  {needs you to open the file and press Run.}
  {opens the file with \menu{Load} when reminded where it is, presses Run, and finds the turtle window with help.}
  {loads the program, presses Run, finds the turtle window and closes it before running again.}
  {Practice the steps with the project starter: \menu{Load}, Run, look, close the window. Put the three steps on a sticky note beside the Mac.}

\smguide{Changes numbers and colors in Python to make it their own (J3.G5)}
  {does not change anything, or changes many things at once and gets lost.}
  {changes a number or a color when you point to the \texttt{CHANGE ME} line.}
  {chooses what to change, changes one thing, presses Run and says what changed in the window (\enquote{I made the house bigger}).}
  {Change one number together and predict first: \enquote{If we make 100 into 150, will the square be bigger or smaller?} Then press Run.}

\smguide{Uses a \texttt{for} loop with an indent (J3.G5, J3.G1)}
  {copies the same line four times instead of using a loop, or the lines inside are not pushed in.}
  {uses the loop from a starter and keeps the indent, but needs help when a line is outside the loop.}
  {changes the number in \texttt{range}, keeps both lines pushed in by the same amount, and says that the pushed-in lines are \enquote{inside the repeat}.}
  {Play \enquote{indent the line-up} from week 7: children who step in are inside the loop. Then compare a Scratch repeat with its Python twin on the cheat sheet.}

\smguide{Fixes a typo error in Python (J3.G6)}
  {sees the red text and asks you to fix it, or deletes lines.}
  {finds the line with the bug when you read the last red line with them, then fixes it.}
  {reads the last red line, goes to that line, compares it with the cheat sheet letter by letter, fixes it and presses Run.}
  {Make a typo on purpose in your own program (\texttt{forwad}) and think aloud while you fix it. Then let the child be the Helper who reads the line letter by letter.}

\smguide{Plans on a plan card and builds it (J3.G2)}
  {starts building without a plan, or the build does not match the plan card.}
  {has a plan card and builds part of it, with reminders to look at it.}
  {fills in the plan card, builds one part at a time, and a partner can check every must-have against the card.}
  {Fill in the plan card together with only the must-haves. Build the first must-have, then check it off on the card.}

\smguide{Explains one part of their code (J3.G5, J3.G6)}
  {says what the project does in general (\enquote{It's a house}).}
  {names some blocks or lines when asked about them.}
  {picks one script or one loop and says what each block or line does, in order, and why it is there.}
  {Practice the 1-minute show-and-tell script with you before the guests arrive. Point at one line and ask \enquote{What does this one do?}}

\smguide{Keeps personal information private and gives credit (J3.G7)}
  {a real full name, school or photo is on the project, or they say they made a library sprite or a starter.}
  {uses a nickname, and says where a starter or idea came from when asked.}
  {uses a nickname only and says, without being asked, \enquote{I started from the postcard starter} or \enquote{The cat is from the Scratch library.} Knows to stop and ask a grown-up if a site asks for their name.}
  {Talk about the week-3, week-4 and week-8 safety boxes again. Change the name on the project together, and add a credit line or say it at the showcase.}

%% ---------------------------------------------------------------------
\section{What to do with the results}
\begin{itemize}[itemsep=2pt]
  \item \textbf{Week 5:} rows with \enquote{Not yet} for most children? Use buffer B1 (Robot Olympics)
        before week 6, or spend the next three warm-ups on those rows. The \enquote{Scratch blocks to
        Python words} row matters most for week 6: if it says \enquote{Not yet}, do the card sort again in
        the week-6 warm-up. One child only? Use the support path in weeks 6--8 and pair them as Helper
        with a patient partner.
  \item \textbf{Week 10:} keep the page with the certificate. Pass it on to whoever teaches main-track
        Level 1 (or keep it for yourself). Rows still at \enquote{Not yet} tell you what to practice in
        the quick review of main-track weeks 1--5.
  \item \textbf{Celebrate growth:} put the week-1, week-5 and week-10 pages side by side with your child
        and count the boxes that moved to the right.
\end{itemize}
